\chapter*{Declaración y uso de IA}
\markboth{Declaración y uso de IA}{Declaración y uso de IA}

Para la redacción y estructuración de esta tesis se empleó Claude, herramienta de
inteligencia artificial generativa de Anthropic, a través de la plataforma Claude Code. Su uso
se enfocó en la redacción y refinamiento de contenidos para los capítulos de introducción,
estado del arte, marco teórico, metodología y resultados. Asimismo, la herramienta asistió en
la revisión gramatical, la reorganización y jerarquización de secciones, y la elaboración tanto
del resumen ejecutivo como de las descripciones para tablas e ilustraciones.
Las instrucciones proporcionadas fueron descriptivas y contextuales, tales como “reescriba
el texto para que sea más claro y conciso”, “resuma la sección en un párrafo de 5 líneas”,
“reorganice la sección para que tenga una estructura más lógica”, “escriba un resumen ejecutivo que capture problema, arquitectura, pruebas experimentales y contribuciones”.
El autor revisó y validó todo el contenido antes de su incorporación. Cada afirmación
técnica se contrastó con datos experimentales y arquitectura efectivamente implementada.
Se verificó coherencia entre secciones, consistencia de notación y línea argumental. Los fragmentos que no correspondían al estilo o nivel de detalle fueron reescritos, y el autor incorporó
datos cuantitativos, referencias bibliográficas y detalles de implementación. Cada capítulo fue
revisado completamente para confirmar que análisis y conclusiones corresponden al criterio
del autor.
El autor revisó y validó todo el contenido de esta memoria. La herramienta se utilizó como
apoyo al trabajo intelectual y no figura como coautora. La responsabilidad académica sobre
los resultados, análisis y conclusiones presentados corresponde íntegramente al autor.